% --- PAQUETES DE IDIOMA Y CODIFICACIÓN ---
\usepackage[utf8]{inputenc}
\usepackage[T1]{fontenc} % Mejor renderizado de fuentes
\usepackage[spanish, es-tabla]{babel} 

% --- CONFIGURACIÓN DE FUENTE (Tipo Arial/Helvética) ---
\usepackage{helvet}
\renewcommand{\familydefault}{\sfdefault} % Establece la fuente Sans como predeterminada

% --- DISEÑO DE PÁGINA ---
\usepackage[margin=2.5cm]{geometry} 
\usepackage{setspace} 
\setstretch{1.15} % Un interlineado cómodo pero no excesivo
\usepackage{graphicx} 
\usepackage{xcolor}   

% --- LISTAS Y TABLAS ---
\usepackage{enumitem} 
\usepackage{booktabs} 
% Definición de colores
\definecolor{codegray}{rgb}{0.5,0.5,0.5}
\definecolor{backcolour}{rgb}{0.98,0.98,0.98}
\definecolor{keyblue}{rgb}{0.13,0.13,1}

\setlist[enumerate,1]{label=\alph*), ref=\alph*)}

\usepackage{listings}

% --- ESTILO DE CAPÍTULOS ---
\usepackage{titlesec}
\titleformat{\chapter}[display]
  {\normalfont\bfseries\huge}
  {\filleft\Large\chaptertitlename\ \thechapter}
  {1ex}
  {\titlerule\vspace{1ex}\filleft}

% --- SOPORTE PARA SUBRAYADO ---
\usepackage[normalem]{ulem} % Permite usar \uline{...} que es mejor que \underline